\chapter{Arquitectura multiagente propuesta}

\noindent Este capítulo presenta la arquitectura multiagente propuesta para transformar datos de mercados secundarios en decisiones de compraventa simuladas. Una arquitectura define cómo se organizan y se comunican los componentes de un sistema para alcanzar un objetivo. En este trabajo, esos componentes reciben datos de los mercados secundarios, preparan la información necesaria y evalúan decisiones sobre una cartera simulada.

La propuesta separa las responsabilidades de adquisición de datos, persistencia, predicción, decisión multiagente y simulación. Esta separación permite distinguir la obtención y preparación de los datos de la decisión posterior. También facilita comprobar cada componente sin confundir una estimación de mercado con una decisión de compra o venta.

Por último, el capítulo define el entorno de cartera, las acciones disponibles, la evolución del estado y la distribución de responsabilidades entre los agentes. Asimismo, presenta las recompensas individuales y cooperativas que orientan el aprendizaje, junto con las restricciones de riesgo que limitan las decisiones. La evaluación se realiza dentro del simulador, sin ejecutar operaciones en los mercados secundarios.

\section{Proceso de decisión simulado}

A partir del planteamiento descrito en el Capítulo 1, la propuesta representa el caso de uso como una tarea de decisión secuencial sobre una cartera simulada, siguiendo la formulación general del aprendizaje por refuerzo \cite{sutton2018reinforcement}. En cada instante, el sistema recibe datos de los mercados secundarios, el estado de la cartera y una estimación supervisada calculada a partir de datos históricos.

La decisión es secuencial porque cada acción modifica las posibilidades disponibles en los instantes posteriores. Una compra reduce el saldo disponible, puede mantener capital bloqueado durante un periodo y genera una posición que condiciona una posible venta posterior. Por ello, el sistema no evalúa cada diferencia de precios de forma aislada, sino como parte de la evolución de la cartera simulada.

La información utilizada para decidir se organiza en tres grupos:
\begin{itemize}
    \item \textbf{Datos de mercado:} artículo, precios publicados, moneda, comisiones y volumen de intercambios.
    \item \textbf{Estado de cartera:} saldo disponible, posiciones abiertas, capital bloqueado y exposición acumulada.
    \item \textbf{Estimación supervisada:} probabilidad obtenida a partir del histórico que orienta la decisión.
\end{itemize}

Con esta información, los agentes coordinan una propuesta de compra, venta o mantenimiento y comprueban si puede realizarse dentro del simulador. Para simular una compra, la publicación empleada para obtener el precio de compra y la publicación utilizada para estimar el precio de venta deben corresponder al mismo artículo y sus precios deben expresarse en una moneda común. Para simular una venta, debe existir una posición del mismo artículo que no esté bloqueada. Las restricciones de cartera pueden impedir que una propuesta se ejecute aunque la diferencia de precios parezca favorable.

El simulador evalúa cada decisión considerando las condiciones de intercambio descritas en el Capítulo 3: conversión de moneda, comisiones, volumen de intercambios, periodo de bloqueo de intercambio y exposición acumulada a un artículo o a una plataforma. El objetivo es incrementar el valor neto de la cartera simulada sin superar los límites de riesgo ni comprometer más capital del permitido.

El flujo comienza cuando los componentes de adquisición recopilan datos de los mercados secundarios y de las fuentes auxiliares. Antes de utilizar los datos extraídos, se normalizan y se conservan junto con su instante de captura para construir un histórico trazable.

Sobre el histórico se calculan las variables necesarias para obtener una estimación supervisada. Una probabilidad alta indica que una posible operación merece pasar a la fase de decisión, pero los agentes todavía deben comprobar el estado de la cartera y las restricciones antes de aprobar una compra simulada. El entorno de cartera combina esa información con el estado de la cartera, genera una observación local para cada agente multiagente y coordina sus acciones dentro del simulador.

\begin{algorithm}
\caption{Flujo de información y decisión simulada}
\begin{algorithmic}[1]
\Require Datos de mercado, estado de cartera y estimación supervisada opcional
\Ensure Cartera actualizada, recompensa y registro trazable
\State Normalizar los datos de mercado y conservar la captura.
\State Consultar el estado de cartera y calcular las restricciones aplicables.
\State Incorporar la estimación supervisada cuando esté disponible.
\State Generar una observación local para cada agente multiagente.
\State Obtener las acciones coordinadas y aplicar las restricciones de riesgo.
\State Simular la transición y actualizar saldo, posiciones y capital bloqueado.
\State Calcular recompensa, métricas y trazabilidad del episodio.
\end{algorithmic}
\end{algorithm}

Cada paso registra la información utilizada, las acciones propuestas y el resultado de la simulación. Este registro permite reconstruir cómo se obtuvo una recomendación y diferenciar la información de partida de las consecuencias que calcula el simulador.

\section{Definición del entorno}

El entorno representa una cartera simulada que evoluciona a partir de datos históricos. Cada \textbf{\glosref{episodio}{episodio}} es una ejecución completa del simulador. Comienza con una cartera inicial, recorre una secuencia ordenada de instantes y termina cuando se han procesado todos los datos del periodo seleccionado. En cada instante, actualiza el saldo disponible, las posiciones abiertas, el capital bloqueado, las fechas de desbloqueo y las métricas de exposición.

La propuesta concreta el proceso de decisión de Markov definido en la Sección~2.3. En este caso, puede expresarse como:

\begin{equation}
\mathcal{M}_{\mathrm{cartera}}=(
\mathcal{S}_{\mathrm{cartera}},
\mathcal{A}_{\mathrm{coordinadas}},
\mathcal{P}_{\mathrm{sim}},
\mathcal{R}_{\mathrm{cartera}},
\gamma).
\label{eq:mdp-cartera}
\end{equation}

En la Ecuación~\ref{eq:mdp-cartera}, $\mathcal{S}_{\mathrm{cartera}}$ describe la situación de la cartera y la información disponible sobre el artículo considerado. $\mathcal{A}_{\mathrm{coordinadas}}$ reúne las acciones propuestas por los agentes. $\mathcal{P}_{\mathrm{sim}}$ determina cómo cambia la cartera después de una acción y $\mathcal{R}_{\mathrm{cartera}}$ evalúa las consecuencias de ese cambio. El factor $\gamma$ indica cuánto peso reciben las recompensas futuras frente a las inmediatas. El objetivo de aprendizaje sigue el criterio general de la Ecuación~\ref{eq:objetivo-rl-general}: obtener una política que maximice el retorno esperado a lo largo del episodio.

Por ejemplo, en un instante concreto el estado puede incluir el saldo disponible, una posición bloqueada y los precios observados para un artículo. Si los agentes aprueban una compra, esa es la acción aplicada. La transición reduce el saldo, abre una posición y registra su fecha de desbloqueo. La recompensa evalúa cómo ha afectado esa decisión al valor neto y al riesgo de la cartera simulada.

El estado global reúne toda la información disponible para el sistema. Una observación local es la parte del estado global que recibe un agente concreto para tomar su decisión. En el entorno se distinguen tres grupos de información:

\begin{itemize}
    \item \textbf{Información del artículo y de las publicaciones:} plataformas de compra y venta, tipo de publicación, precio de compra normalizado, precio de venta previsto, retorno previsto y cantidad disponible.
    \item \textbf{Estado de cartera:} saldo disponible, saldo previsto después de la compra, capital bloqueado, exposición asociada al artículo y posiciones que ya pueden venderse o que siguen bloqueadas.
    \item \textbf{Estimación supervisada:} probabilidad disponible para el artículo y un indicador que señala si esa estimación puede utilizarse en el episodio.
\end{itemize}

La estimación supervisada no obliga al sistema a comprar. Una probabilidad alta solo aporta una señal adicional sobre el artículo. Si la cartera supera alguno de sus límites o el capital disponible no permite una compra, la operación puede rechazarse o posponerse dentro del simulador.

El espacio de acciones es discreto y se limita a ignorar o marcar un artículo, mantener la cartera, comprar una unidad, vender una posición abierta, aprobar una propuesta o rechazarla. El entorno calcula antes de cada paso qué acciones pueden ejecutarse. Por ejemplo, no permite vender si no existe una posición desbloqueada del artículo, ni comprar si la operación supera los límites configurados.

La exploración y el entrenamiento se realizan únicamente dentro del simulador y con datos históricos. Las acciones modifican una cartera simulada, no una cuenta ni artículos reales.

\section{Agentes de decisión}

La decisión simulada se distribuye entre tres agentes \acroref{MARL}{MARL} con responsabilidades diferenciadas. Cada agente recibe una observación local adaptada a su función y propone una acción dentro del entorno de simulación:

\begin{itemize}
    \item \textbf{Scout:} revisa la plataforma, el tipo de publicación, el precio de compra, el retorno previsto, la disponibilidad y la probabilidad supervisada. Su función es identificar los artículos que merecen atención.
    \item \textbf{Trader:} utiliza la señal de Scout, la información de compra y venta prevista y las posiciones del mismo artículo que ya existen en la cartera para proponer la acción operativa principal.
    \item \textbf{Portfolio:} comprueba el saldo, el capital bloqueado, la exposición prevista, las posiciones que pueden venderse y las restricciones de riesgo antes de validar una propuesta.
\end{itemize}

Las acciones disponibles para cada agente son las siguientes:
\begin{itemize}
    \item \textbf{Scout:} \textit{ignorar} o \textit{marcar} un artículo.
    \item \textbf{Trader:} \textit{mantener} la cartera, \textit{comprar} una unidad o \textit{vender} una posición desbloqueada del artículo.
    \item \textbf{Portfolio:} \textit{rechazar} o \textit{aprobar} la propuesta.
\end{itemize}

\begin{algorithm}[H]
\caption{Decisión cooperativa}
\label{alg:decision-cooperativa}
\begin{algorithmic}[1]
\Require Artículo candidato $c_t$ y observaciones locales
\Ensure Transición de cartera y recompensa compartida $r_t$
\State $a_t^{S} \gets \pi_S(o_t^{S})$
\State $a_t^{T} \gets \pi_T(o_t^{T})$
\State $a_t^{P} \gets \pi_P(o_t^{P})$
\If{$a_t^{S}=\textit{marcar}$ \textbf{y} $a_t^{T}=\textit{comprar}$ \textbf{y} $a_t^{P}=\textit{aprobar}$ \textbf{y} la compra cumple los límites de riesgo}
  \State Abrir una posición simulada y bloquear su capital
\ElsIf{$a_t^{T}=\textit{vender}$ \textbf{y} $a_t^{P}=\textit{aprobar}$ \textbf{y} existe una posición desbloqueada de $c_t$}
  \State Cerrar una posición simulada y abonar su importe neto
\Else
  \State Mantener el saldo y registrar el motivo de no ejecución
\EndIf
\State Actualizar el saldo, las posiciones, la exposición y las métricas de riesgo.
\State Calcular $r_t$ y devolver la misma recompensa a los tres agentes.
\end{algorithmic}
\end{algorithm}

En el Algoritmo~\ref{alg:decision-cooperativa}, $c_t$ representa el artículo considerado en el instante $t$. Las funciones $\pi_S$, $\pi_T$ y $\pi_P$ son las políticas de Scout, Trader y Portfolio. Cada política recibe su observación local, representada por $o_t^S$, $o_t^T$ u $o_t^P$, y produce una acción, representada por $a_t^S$, $a_t^T$ o $a_t^P$.

Una compra simulada requiere que Scout marque el artículo, Trader solicite la compra y Portfolio la apruebe. Una venta simulada requiere la solicitud de Trader, la aprobación de Portfolio y una posición desbloqueada del mismo artículo. Scout no emite una marca nueva para vender porque la posición ya forma parte de la cartera. La recompensa compartida que recibe cada agente se explica en la Sección~4.4.

\section{Función de recompensa}

La recompensa debe reflejar el objetivo del sistema: mejorar el valor neto de la cartera sin ignorar riesgo ni fricción. Una recompensa basada solo en comprar barato o en detectar subida de precio sería incompleta, porque no tendría en cuenta si la operación puede cerrarse, si el capital queda bloqueado o si el spread desaparece antes de poder vender.

\subsection{Recompensa individual}

Las señales individuales pueden ayudar a entrenar comportamientos especializados. Scout puede recibir recompensa por identificar oportunidades que más adelante resultan útiles. Trader puede recibirla por ejecutar acciones con beneficio neto y Portfolio por evitar sobreexposición, operaciones ilíquidas o capital bloqueado excesivo. Estas recompensas no sustituyen al objetivo global, pero pueden guiar el aprendizaje de cada rol.

Como ampliación de la arquitectura, la recompensa de cada rol puede combinar una parte compartida, vinculada a la evolución de la cartera, y una parte individual, vinculada a su responsabilidad concreta:

\begin{equation}
r_t^i=\lambda r_t^{\mathrm{cartera}}+(1-\lambda)r_{t,\mathrm{rol}}^i,
\qquad 0\leq\lambda\leq1.
\label{eq:recompensa-hibrida}
\end{equation}

En la Ecuación~\ref{eq:recompensa-hibrida}, $r_t^{\mathrm{cartera}}$ es la recompensa común, $r_{t,\mathrm{rol}}^i$ mide la contribución del agente $i$ y $\lambda$ determina el peso de cada parte. La parte compartida debe tener un peso mayor para que todos los agentes mantengan como prioridad el resultado global de la cartera. Las ejecuciones descritas en este trabajo utilizan la recompensa compartida de la sección siguiente. La recompensa individual queda definida como una extensión que deberá evaluarse con datos suficientes.

\subsection{Recompensa cooperativa}

La recompensa cooperativa se vincula al valor total de la cartera simulada. Considera beneficio realizado, retorno de la operación ejecutada, inactividad ante una oportunidad permitida, restricciones de riesgo, \textit{drawdown}, capital bloqueado y volatilidad. La expresión implementada es:

\[
r_t = \frac{\Delta \Pi_t}{100}
+ \rho_t\mathbb{I}_{\mathrm{operacion}}
- 0{,}01\mathbb{I}_{\mathrm{oportunidad\ ignorada}}
- 0{,}05v_t
- 0{,}50DD_t
- 0{,}05B_t
- 0{,}10\sigma_t\mathbb{I}_{\mathrm{compra}}.
\]

Aquí, $\Delta \Pi_t$ es el cambio en beneficio realizado en euros, $\rho_t$ es el retorno de la compra o venta ejecutada, $v_t$ cuenta las restricciones incumplidas, $DD_t$ es el \textit{drawdown} relativo, $B_t$ es la fracción de cartera bloqueada y $\sigma_t$ representa la volatilidad disponible para el candidato. El divisor 100 evita que una ganancia expresada en euros domine al resto de términos. La misma recompensa se asigna a Scout, Trader y Portfolio, es decir, $r_t^S=r_t^T=r_t^P=r_t$, para que los tres roles persigan el estado global de la cartera.

\section{Restricciones de riesgo}

Las restricciones de riesgo se introducen antes de escalar el entrenamiento. Las más importantes son: no usar capital inexistente, limitar exposición por activo o plataforma, rechazar operaciones con un volumen de intercambios insuficiente, respetar el \textit{trade hold}, controlar el capital bloqueado y separar claramente simulación de operación. Estas restricciones de riesgo son sencillas, pero evitan que los agentes aprendan políticas imposibles de ejecutar en el mercado.

La salida del sistema debe ser siempre auditable. Cada predicción, acción o decisión simulada debe poder relacionarse con su fuente de datos, versión de modelo, episodio y configuración de simulación. Esta trazabilidad es necesaria para interpretar resultados y para distinguir entre error de datos, error de modelo y mala política de decisión.

\section{Síntesis de la arquitectura propuesta}

La arquitectura se organiza en cinco capas: adquisición, persistencia, predicción, decisión multiagente y simulación. Cada capa tiene una responsabilidad concreta y se comunica con la siguiente mediante datos normalizados. Esta decisión de diseño reduce el acoplamiento y permite validar cada parte por separado, una práctica habitual en entornos de trading reproducibles y backtesting \cite{liu2021finrl}.

\begin{figure}[h]
\centering
\fbox{
\begin{minipage}{0.92\textwidth}
\centering
\begin{tabularx}{0.95\textwidth}{>{\bfseries}l X}
1. Adquisición & Steam, \acroref{BUFF}{BUFF}, SteamDT, \acroref{OCR}{OCR} e importaciones manuales capturan observaciones. \\
2. Persistencia & Postgres/Supabase guarda artículos, históricos, snapshots, predicciones y decisiones simuladas. \\
3. Predicción & El modelo supervisado genera probabilidades calibradas y señales de mercado. \\
4. Decisión \acroref{MARL}{MARL} & Scout, Trader y Portfolio deciden dentro del entorno PettingZoo/\acroref{RLlib}{RLlib}. \\
5. Simulación & La cartera aplica comisiones, divisas, \textit{trade hold}, riesgo y métricas. \\
\end{tabularx}
\end{minipage}
}
\caption{Capas del sistema.}
\end{figure}

La decisión final siempre queda registrada como simulada. Esto permite evaluar el sistema con métricas de cartera sin exponer capital durante la investigación. La arquitectura puede generar recomendaciones, episodios y métricas, pero no envía órdenes a Steam, \acroref{BUFF}{BUFF} ni ninguna otra plataforma.
